% GENERATED FILE. Edit canonical lessons, then run npm run generate:books.
% canonical-source-hash: fnv1a64:445a87b7a04b15a2
% canonical-lessons: JA-C530-kanari, JA-C530-sukunakutomo, JA-C530-hobo, JA-C530-oyoso, JA-C530-mattaku

\chapter{At Least, About, Almost}
\label{ch:ja-at-least-about-almost}
\hlchaptermodality{530}

\begin{chapteropening}
\textbf{By the end of this chapter:} \emph{I can say roughly how long and how much.}

Complete the last lesson of chapter 530: I can say roughly how long and how much.
\end{chapteropening}

\section[\texorpdfstring{\ja{かなり} (kanari)}{かなり (kanari)}]{\ja{かなり} (kanari) — fairly; quite}

\label{lesson:JA-C530-kanari}



\begin{quote}
Before the new one: say the Japanese for would sooner, then the Japanese for as one would expect.
\end{quote}

\subsection*{You'll want to know: \ja{かなり}}
\textbf{\ja{かなり}} — \emph{kanari} — \textquotedblleft{}fairly; quite\textquotedblright{}.

\begin{quote}
\textbf{\ja{もし} \ja{あめが} \ja{ふったら、みちは} \ja{かなり} \ja{こむでしょう。}} — \emph{moshi ame ga futtara, michi wa kanari komu deshō} — \textquotedblleft{}If it rains, the roads will be fairly crowded.\textquotedblright{}
\end{quote}

\begin{cousinweb}
\textbf{\ja{かなり}} (\emph{kanari}) is usually counted a native Japanese word (wago), inherited rather than borrowed from Chinese, though that is not certain: JMdict marks it as usually written in kana and lists only rare kanji spellings, and no dictionary we use records its history, so where it comes from is not settled (JMdict).
\end{cousinweb}

\begin{grammarlens}[title={What you've built}]
One more word for saying what happens if.
\end{grammarlens}

\subsection*{Your turn}
\begin{itemize}
\raggedright
  \item \emph{Say it:} \emph{kanari}
  \item \emph{Say it:} \emph{kanari}, once more
  \item \emph{Say it:} \emph{moshi ame ga futtara, michi wa kanari komu deshō} — If it rains, the roads will be fairly crowded.
  \item \emph{Recall:} say the Japanese for would sooner, then the Japanese for as one would expect, then say \emph{kanari} again
\end{itemize}

\begin{tcolorbox}[breakable,colback=teal!4,colframe=teal!35!black,title={Before you move on}]
What does \ja{かなり} mean? (\textquotedblleft{}Fairly; quite\textquotedblright{}.) Say it once more.
\end{tcolorbox}
\section[\texorpdfstring{\ja{すくなくとも} (sukunakutomo)}{すくなくとも (sukunakutomo)}]{\ja{すくなくとも} (sukunakutomo) — at the very least; no fewer than}

\label{lesson:JA-C530-sukunakutomo}



\begin{quote}
Before the new one: say the Japanese for as one would expect, then the Japanese for fairly.
\end{quote}

\subsection*{You'll want to know: \ja{すくなくとも}}
\textbf{\ja{すくなくとも}} — \emph{sukunakutomo} — \textquotedblleft{}at the very least; no fewer than\textquotedblright{}.

\begin{quote}
\textbf{\ja{もし} \ja{くるまで} \ja{いったら、すくなくとも} \ja{いちじかん} \ja{かかるでしょう。}} — \emph{moshi kuruma de ittara, sukunakutomo ichijikan kakaru deshō} — \textquotedblleft{}If you go by car, it will take at least an hour.\textquotedblright{}
\end{quote}

\begin{cousinweb}
\textbf{\ja{すくなくとも}} (\emph{sukunakutomo}) is a native Japanese word (wago), inherited rather than borrowed from Chinese: built on \emph{sukunai}, \textquotedblleft{}few\textquotedblright{}, with the ending \emph{tomo}, \textquotedblleft{}even\textquotedblright{}; its kanji was attached later as a kun'yomi (JMdict, KANJIDIC2).
\end{cousinweb}

\begin{grammarlens}[title={What you've built}]
One more word for saying what happens if.
\end{grammarlens}

\subsection*{Your turn}
\begin{itemize}
\raggedright
  \item \emph{Say it:} \emph{sukunakutomo}
  \item \emph{Say it:} \emph{sukunakutomo}, once more
  \item \emph{Say it:} \emph{moshi kuruma de ittara, sukunakutomo ichijikan kakaru deshō} — If you go by car, it will take at least an hour.
  \item \emph{Recall:} say the Japanese for as one would expect, then the Japanese for fairly, then say \emph{sukunakutomo} again
\end{itemize}

\begin{tcolorbox}[breakable,colback=teal!4,colframe=teal!35!black,title={Before you move on}]
What does \ja{すくなくとも} mean? (\textquotedblleft{}At the very least; no fewer than\textquotedblright{}.) Say it once more.
\end{tcolorbox}
\section[\texorpdfstring{\ja{ほぼ} (hobo)}{ほぼ (hobo)}]{\ja{ほぼ} (hobo) — nearly; almost}

\label{lesson:JA-C530-hobo}



\begin{quote}
Before the new one: say the Japanese for fairly, then the Japanese for at the very least.
\end{quote}

\subsection*{You'll want to know: \ja{ほぼ}}
\textbf{\ja{ほぼ}} — \emph{hobo} — \textquotedblleft{}nearly; almost\textquotedblright{}.

\begin{quote}
\textbf{\ja{もし} \ja{ごじに} \ja{でたら、ちちは} \ja{ほぼ} \ja{まにあうでしょう。}} — \emph{moshi goji ni detara, chichi wa hobo maniau deshō} — \textquotedblleft{}If he leaves at five, my father will make it, more or less.\textquotedblright{}
\end{quote}

\begin{cousinweb}
\textbf{\ja{ほぼ}} (\emph{hobo}) is a native Japanese word (wago), inherited rather than borrowed from Chinese, and usually written in kana; its older origin is not settled (JMdict).
\end{cousinweb}

\begin{grammarlens}[title={What you've built}]
One more word for saying what happens if.
\end{grammarlens}

\subsection*{Your turn}
\begin{itemize}
\raggedright
  \item \emph{Say it:} \emph{hobo}
  \item \emph{Say it:} \emph{hobo}, once more
  \item \emph{Say it:} \emph{moshi goji ni detara, chichi wa hobo maniau deshō} — If he leaves at five, my father will make it, more or less.
  \item \emph{Recall:} say the Japanese for fairly, then the Japanese for at the very least, then say \emph{hobo} again
\end{itemize}

\begin{tcolorbox}[breakable,colback=teal!4,colframe=teal!35!black,title={Before you move on}]
What does \ja{ほぼ} mean? (\textquotedblleft{}Nearly; almost\textquotedblright{}.) Say it once more.
\end{tcolorbox}
\section[\texorpdfstring{\ja{およそ} (oyoso)}{およそ (oyoso)}]{\ja{およそ} (oyoso) — about; roughly}

\label{lesson:JA-C530-oyoso}



\begin{quote}
Before the new one: say the Japanese for at the very least, then the Japanese for nearly.
\end{quote}

\subsection*{You'll want to know: \ja{およそ}}
\textbf{\ja{およそ}} — \emph{oyoso} — \textquotedblleft{}about; roughly\textquotedblright{}.

\begin{quote}
\textbf{\ja{でんしゃで} \ja{いったら、およそ} \ja{にじかん} \ja{かかるでしょう。}} — \emph{densha de ittara, oyoso nijikan kakaru deshō} — \textquotedblleft{}If you go by train, it will take about two hours.\textquotedblright{}
\end{quote}

\begin{cousinweb}
\textbf{\ja{およそ}} (\emph{oyoso}) is a native Japanese word (wago), inherited rather than borrowed from Chinese, and usually written in kana; its older origin is not settled (JMdict).
\end{cousinweb}

\begin{grammarlens}[title={What you've built}]
One more word for saying what happens if.
\end{grammarlens}

\subsection*{Your turn}
\begin{itemize}
\raggedright
  \item \emph{Say it:} \emph{oyoso}
  \item \emph{Say it:} \emph{oyoso}, once more
  \item \emph{Say it:} \emph{densha de ittara, oyoso nijikan kakaru deshō} — If you go by train, it will take about two hours.
  \item \emph{Recall:} say the Japanese for at the very least, then the Japanese for nearly, then say \emph{oyoso} again
\end{itemize}

\begin{tcolorbox}[breakable,colback=teal!4,colframe=teal!35!black,title={Before you move on}]
What does \ja{およそ} mean? (\textquotedblleft{}About; roughly\textquotedblright{}.) Say it once more.
\end{tcolorbox}
\section[\texorpdfstring{\ja{まったく} (mattaku)}{まったく (mattaku)}]{\ja{まったく} (mattaku) — (not) at all; entirely}

\label{lesson:JA-C530-mattaku}



\begin{quote}
Before the new one: say the Japanese for nearly, then the Japanese for about.
\end{quote}

\subsection*{You'll want to know: \ja{まったく}}
\textbf{\ja{まったく}} — \emph{mattaku} — \textquotedblleft{}(not) at all; entirely\textquotedblright{}.

\begin{quote}
\textbf{\ja{もし} \ja{ちずを} \ja{わすれたら、みちが} \ja{まったく} \ja{わかりません。}} — \emph{moshi chizu o wasuretara, michi ga mattaku wakarimasen} — \textquotedblleft{}If you forget the map, you will not know the way at all.\textquotedblright{}
\end{quote}

\begin{cousinweb}
\textbf{\ja{まったく}} (\emph{mattaku}) is a native Japanese word (wago), inherited rather than borrowed from Chinese, and usually written in kana; its older origin is not settled (JMdict).
\end{cousinweb}

\begin{grammarlens}[title={What you've built}]
One more word for saying what happens if.
\end{grammarlens}

\subsection*{Your turn}
\begin{itemize}
\raggedright
  \item \emph{Say it:} \emph{mattaku}
  \item \emph{Say it:} \emph{mattaku}, once more
  \item \emph{Say it:} \emph{moshi chizu o wasuretara, michi ga mattaku wakarimasen} — If you forget the map, you will not know the way at all.
  \item \emph{Recall:} say the Japanese for nearly, then the Japanese for about, then say \emph{mattaku} again
\end{itemize}

\begin{tcolorbox}[breakable,colback=teal!4,colframe=teal!35!black,title={Before you move on}]
What does \ja{まったく} mean? (\textquotedblleft{}(not) at all; entirely\textquotedblright{}.) Say it once more.
\end{tcolorbox}
